%------------------------------------------------------------------------------%
% \chapter{Revisão}
%------------------------------------------------------------------------------%

%------------------------------------------------------------------------------%
\begin{exercicios}{Expoentes e Radicais}
    
    \begin{exercicio}
        Calcule os valores das expressões.
        \begin{mathlist}[2]
            \item 27^{2/3}
            \item 8^{-4/3}
            \item \left(\frac{1}{\sqrt{3}}\right)^0
            \item (7^{1/2})^4
            \item \left[\left(\frac{1}{8}\right)^{1/3}\right]^{-2}
            \item \left[\left(- \frac{1}{3}\right)^2\right]^{-3}
            \item \left(\frac{7^{-5} \, 7^2}{7^{-2}}\right)^{-1}
            \item (125^{2/3})^{-1/2}
            \item \sqrt[3]{2^6}
            \item \frac{\sqrt{32}}{\sqrt{8}}
            \item \sqrt[3]{\frac{-8}{27}}
            \item \frac{16^{5/8} \, 16^{1/2}}{16^{7/8}}
            \item \frac{6^{2,5} \, 6^{-1,9}}{6^{-1,4}}
        \end{mathlist}	
        \begin{answers}
            \begin{mathlist}[3]
                \item 9
                \item \frac{1}{16}
                \item 1
                \item 49
                \item 4
                \item 729
                \item 7
                \item \frac{1}{5}
                \item 4
                \item 2
                \item \frac{2 \sqrt[3]{-1}}{3}
                \item 2
                \item 36
            \end{mathlist}
        \end{answers}
    \end{exercicio}

\todo{    
    \begin{exercicio}
        Copiar exercícios 1-6,7-12,13-18, pg 321 livro do Tan
    \end{exercicio}
}
    
    \begin{exercicio}
        Determine se as afirmações são verdadeiras ou falsas. Justifique sua decisão.	
        \begin{mathlist}[2]
            \item x^4 + 2x^4 = 3x^4
            \item 3^2 \, 2^2 = 6^2
            \item \frac{2^{4x}}{1^{3x}} = 2^{4x-3x}
            \item (2^2 \, 3^2)^2 = 6^4
            \item \frac{4^{3/2}}{2^4} = \frac{1}{2}
            \item (1,2^{1/2})^{-1/2} = 1		
        \end{mathlist}	
    \end{exercicio}
    
    \begin{exercicio}
        Reescreva as expressões utilizando apenas expoentes positivos.
        \begin{mathlist}[2]
            \item (xy)^{-2}
            \item \frac{x^{-1/3}}{x^{1/2}}
            \item \sqrt{x^{-1}} \, \sqrt{9x^{-3}}
            \item 12^0 (s + t)^{-3}
        \end{mathlist}			
    \end{exercicio}
    
    \begin{exercicio}
        Simplifique as expressões. (Assuma que $x$, $y$, $r$, $s$ e $t$ são positivos.)
        \begin{mathlist}[2]
            \item \frac{x^{7/3}}{x^{-2}}
            \item (49x^{-2})^{1/2}
            \item (x^2y^{-3})(x^{-5}y^3)
            \item \frac{5x^6 y^3}{2x^2 y^7}
            \item \frac{x^{3/4}}{x^{-1/4}}
            \item \left(\frac{x^3}{-27y^{-6}}\right)^{-2/3}
            \item \left(\frac{x^{-3}}{y^{-2}}\right)^2 \left(\frac{y}{x}\right)^4
            \item \frac{(r^n)^4}{r^{5-2n}}
            \item \sqrt[3]{x^{-2}} \, \sqrt{4x^5}
            \item \sqrt{81x^6 y^{-4}}
            \item \sqrt[3]{27r^6} \, \sqrt{s^2 t^4}
        \end{mathlist}	
    \end{exercicio}
    
    \begin{exercicio}
        Utilize o fato de que $2^{1/2} \approx 1,414$ e $3^{1/2} \approx 1,732$ para calcular
        as expressões dadas sem o auxílio de uma calculadora.	
        \begin{mathlist}[2]
            \item 2^{3/2}
            \item 8^{1/2}
        \end{mathlist}	
    \end{exercicio}
    
    \begin{exercicio}
        Utilize o fato de que $10^{1/2} \approx 3,162$ e $10^{1/3} \approx 2,154$	para calcular
        as expressões dadas sem o auxílio de uma calculadora.	
        \begin{mathlist}[3]
            \item 10^{3/2}
            \item 10^{2,5}
            \item (0,0001)^{-1/3}
        \end{mathlist}	
    \end{exercicio}
    
    \begin{exercicio}
        Racionalize o denominador das expressões.
        \begin{mathlist}[3]
            \item \frac{3}{2 \sqrt{x}}
            \item \frac{5x^2}{\sqrt{3x}}
            \item \sqrt{\frac{2x}{y}}
        \end{mathlist}	
    \end{exercicio}
    
    \begin{exercicio}
        Racionalize o denominador das expressões.
        \begin{mathlist}[2]
            \item \frac{2\sqrt{x}}{3}
            \item \frac{\sqrt[3]{x}}{24}
            \item \sqrt{\frac{2y}{x}}
            \item \frac{\sqrt[3]{x^2 z}}{y}
        \end{mathlist}
    \end{exercicio}
    
\end{exercicios}

%------------------------------------------------------------------------------%
